% !TeX root = ../papers/nc-p3.tex
\section{有限な低階数軌跡の核グラフとその標準偏極}
\label{sec:octic-kernel-canonical-model}
\begin{theorem}\label{thm:octic-kernel-canonical-model}
定理\ref{thm:octic-kernel-image-degree}の有限な階数2以下軌跡の場合、核グラフの主成分 $G$ を取る。
その入力・出力超平面束を加法記法で $H,L$ と書く。$G$ はGorenstein整曲面であり、
標準束は非常に豊富、$K_G^2=68,\ h^1(\mathcal O_G)=0,\ p_g(G)=15$ である。
\[
             \tau=K_G-H-L\quad\text{は非自明な二捩れである。}
\]
二つの射影はともに完全な四切断で与えられ、$H+L$ の切断は十六個である。
$\tau$ の連結エタール二重被覆も非常に豊富な標準束を持ち、$K^2=136,\ h^1=0,\ p_g=31$ となる。
これらの主張に $X$ の被約性を要しない。
\end{theorem}
\begin{proof}[点窓の吹上げと実際の核グラフ]
$B=\operatorname{Bl}_X\mathbb P^3$、例外Cartier因子を $E$ とする。
六次余因子は $\mathcal I_X^2$ を生成するので、その有理写像のグラフは二乗イデアルの吹上げ、
即ち $B$ 自体である。射影先の超平面束は $6H-2E$ である。$D$ の稠密開ではこの写像が
核の二次Veroneseなので、その真変換は $G$ に同型で $2L=(6H-2E)|_G$ である。
$G$ 自身は $B$ 上のCartier因子 $8H-3E$ である。実際 $I_X$ は局所三元正則列で
$E$ は $X$ 上の射影平面束、$d$ の三次初期形式は各台点の相対射影平面で前節の非零三次式 $\det S_1$ となる。
従って $3E$ を引いた後は垂直な因子成分を持たない。正則列の吹上げは各チャートで
他二生成元を選んだ生成元の倍と置く二式の完全交叉である。よってこのCartier因子は純で、
稠密開で被約かつ既約なので真変換に一致する。吹上げの双対化層は $-4H+2E$ であり、
局所的にも三生成元の比を加える二式に随伴を適用すればよい。
従って $K_G=(4H-E)|_G$ で、上の二乗式は $2\tau=0$ を与える。

$4H-E$ は $B$ 上で非常に豊富である。実際この写像は指定辺の一次六行四列行列 $A(x)$ の
十五の最大行列式であり、像は $\operatorname{Gr}(4,6)$ に入る。ある四空間 $U$ で
$\operatorname{im} A(x)\subset U$ となる入力ベクトル空間は高々一次元である。
二次元以上なら入力直線の上で $A$ は四行四列の行列になり、$X$ の制限は行列式一つで切られる。
これは全直線または長さ4となり、直線切断上界に反する。従ってグラフの像上では
$x\mapsto A(x)\bmod U$ の核が入力点を正則に復元する（単元三次小行列式による消去）。
この写像は $B$ を閉埋込みにする。実際 $B$ は辺の最大行列式（$\mathcal I_X$ の生成元）のグラフであり、
上記の逆射とGrassmann多様体のPlücker埋込みを用いればよい。

$H^0(G,(3H-E)|_G)=0$ である。吹上げ上では三次評価の単射性から切断がない。
制限完全列の左端 $-5H+2E$ の第一コホモロジーも零である。吹上げ射で
$\mathcal O_B(jE)$ の直像は $\mathcal O$、高次直像は $0$（$0\le j\le2$）である。
$j=0$ は相対Serre消滅から $\mathcal O_B(-mE)$ の列を降り、正則列イデアルの次数商
$I^m/I^{m+1}$ が例外束の $\mathcal O_E(m)$ の直像に等しいことを使う。
$j=1,2$ は例外平面上の $\mathcal O(-1),\mathcal O(-2)$ の全コホモロジー消滅による。
射影空間の $H^1(\mathcal O(-5))=0$ から所望となる。出力には非零切断があるため $L\ne(3H-E)|_G$、すなわち $\tau\ne0$ である。
\end{proof}

\begin{proof}[リンケージと二重被覆]
曲面の完全交叉入射を $Y$、$X$ 上で $y\cdot v=0$ を課す射影平面束を $V$ とする。
前節の重複度の証明により $V$ を除いた残余は主成分 $G$ そのものである。
完全交叉リンケージより残余はCohen--Macaulayで、一般点で被約かつ他の成分がないためである。
$\omega_Y=\mathcal O_Y(4,0)$ であり、前節と同じ双対完全列の左端は $V$ の各Artin局所成分で
$\mathcal O_{\mathbb P^2_X}(-3)$ である。Koszul複体より $H^1(Y,\mathcal O_Y)=0$、
$H^2(Y,\mathcal O_Y)$ の双対は入力全四次式の35次元空間である。
左端からの第二コホモロジーの射の双対は $X$ への四次評価であり全射である。
従って $h^1(G)=0$, $p_g=35-20=15$ を得る。
$(0,1),(1,1)$ でひねると左端の平面上の次数は $-2$、従って全コホモロジーが消え、
周囲の四または十六切断がそのまま $G$ の完全切断である。$(1,0)$ では左端の $H^0,H^1$ が零で同様である。
$H^2=8,H L=24,L^2=12$ は主成分と垂直平面の交叉数なので $K_G^2=(H+L)^2=68$。
$\tau$ の二重被覆は連結で、標準切断は不変部分 $H^0(K_G)$ と反不変部分 $H^0(H+L)$ の和である。
従って $p_g=31$、エタール引戻しで $K^2=136$ となる。Euler数は32である。
例えばねじれ線束の有理Chern指標は1なのでRiemann--Rochで元の二倍となる。
従って $h^1=0$。不変部分で基底上の点と接方向を分離し、基点なしの反不変部分で二枚を分離するので標準束は非常に豊富である。
\end{proof}

\begin{corollary}\label{cor:octic-canonical-marked-matrix}
この範囲の反対行列の入力基底変更と出力合同による同値は、$(G,H,\tau)$ の同型と同値である。
一つの抽象曲面 $G$ に対応する行列類は有限である。
\end{corollary}
\begin{proof}
$L=K_G-H+\tau$ を復元し、二つの完全線形系でグラフを戻す。
グラフで消える $(2,1)$ の式は元の四行だけである。実際 $\mathcal I_{G/Y}(2,1)$ の $H^0$ は
垂直平面上の次数 $-2$ により零で、$Y$ のこの次数の式は完全交叉の定義四式に限られる。
同じ行空間の二つの対称表示が $N,CN$ なら $CN=NC^{\mathsf t}$。$N(x)$ の値は全対称行列を張るので
$C$ は定数スカラーである。逆に行列の同値はグラフとその二つの偏極を保つ。
固定した $G$ で $H^2=8,K_G H=32$、$K_G$ は豊富である。Hodge指数定理により
これらの条件を満たす数値格子点は有限である（$K_G$ の直交補上の負定値性）。
整曲面でも解消上に引き戻して同じ指数定理を使える。$h^1(\mathcal O_G)=0$ とNéron--Severiの
有限生成性から各数値類の線束および二捩れ線束の選択は有限である。
これは反対行列についての同値で、中央辺を伴う四窓の判定には別途その窓の条件を使う。
\end{proof}

\path{computations/verify_octic_kernel_canonical_model.py} は完全交叉と垂直平面束のEuler数から
二変数Hilbert多項式、標準自己交叉数および二重被覆の数値を確認する。
